% =========================================================================
\section{Input admissibility: withhold supervision, not images}
\label{sec:ch4-admissibility}
% =========================================================================
The policy of \cref{sec:ch4-conditional-adaptation} acts on every buffered sample once its label arrives. This section asks whether some samples should never reach it, the first condition of the adaptation protocol and the only one decidable at $t$, from the input alone, before any label has arrived. The originating publication~\cite{Piechocki2025:i63677} proposed an anomaly score on the forecasting embedding as that screen. The evidence assembled here, three detectors and the policy run on four fault streams, reverses the premise. A score on the pooled embedding measures forecasting difficulty and is blind to sensing faults; a per-position score on an early encoder stage detects them; and, once they are detected, withholding them turns out to be the wrong response for every fault that persists. A corrupted frame whose label is still valid is a training example, and the policy learns to see through dirt or around an occluder within weeks if it is allowed to. What a screen may legitimately withhold is a transient event, which the policy cannot learn from, and a frame whose label is wrong, which this section does not test. The remainder specifies the detectors, what each one detects, and what the policy does with and without them.

\subsection{Method}
\label{sec:ch4-admissibility-method}
Not every departure from the training distribution is a fault, and not every fault is unlearnable. Dirt, insects and birds on the dome, water on it, and acquisition faults in the camera all change the frame without changing the irradiance it must predict. The methods below were developed in sequence, each pre-registered with its pass criteria before it was scored, and the sequence is reported as it happened, including the designs that failed.

\paragraph{Data and detector fitting.} All detectors are fitted on the Folsom record in the second partition of \cref{tab:ch2-shared-assets}. The training and validation partitions, already used to fit the forecaster, additionally serve to fit the detectors, and the held-out 2016 year is replayed as a stream of previously unseen samples. Every fifth frame of 2014--2015 is used, split by whole days: 458 days (\num{62251} frames) fit each detector and 115 held-out days (\num{15466} frames) set its threshold, always at the 99th percentile of the detector's own calibration scores, a nominal false-alarm rate of \SI{1}{\percent} following \cref{sec:ch2-ood}.

\paragraph{Embedding-space scoring.} The reused ResNet50 encoder maps each frame to $z_t = g(x_t;\theta_b) \in \mathbb{R}^{\num{2048}}$. An Isolation Forest~\cite{articleLiu2012} with 100 estimators is fitted to nominal training embeddings and applied unchanged; a Mahalanobis distance and a $k$-nearest-neighbour distance on the same embeddings, and two model-free references on raw pixels, a dark-pixel fraction and a Laplacian-variance blur score, are compared with it. Scoring traverses the fitted trees without a second encoder pass, so its added work depends on the forest rather than on image resolution. This is the screen of the originating publication, retained here because what it fails to detect is instructive.

\paragraph{Per-position scoring.} The camera never moves, so every position in the frame has its own notion of normal. The second detector reads the encoder's second stage, a $512 \times 16 \times 16$ tensor at the $128 \times 128$ input, so one cell per $8 \times 8$ pixel block, pools each cell with its $3 \times 3$ neighbourhood and projects it to 128 dimensions with a fixed random matrix. For each of the 172 cells whose block lies inside the circular field of view, a memory bank holds that cell's features from \num{2000} nominal frames; a cell's score is its nearest-neighbour distance to its own bank, divided by the 95th percentile of the same distance on \num{1000} other nominal frames. The design follows the memory-bank detectors of industrial inspection~\cite{roth2022towards}, transposed to a fixed camera. The bank occupies \SI{176}{\mega\byte} in \gls{fp32}; scoring latency on the target platform is not measured. Because the second stage is below the encoder stages that the damped action of \cref{eq:ch4-damped-encoder-update} updates, the detector is unaffected by adaptation and needs no refitting as the policy runs.

\paragraph{The exclusion flag.} A frame score is some statistic of its 172 cell scores, and the choice of statistic decides which faults are visible. The maximum over all cells, the first pre-registered design, is dominated by the outermost ring of cells, where the low sun and its reflections sit: \SI{95}{\percent} of the frames it flags on the clean 2016 stream are blamed on that ring, and three quarters of them fall at solar elevations below \ang{20} or within two cells of the projected sun. The design retained, the third pre-registered, uses two statistics chosen for two footprints. The \emph{global} score $G$ is the median over all cells and responds to frame-wide departures, an acquisition fault or water over the dome; the \emph{core} score $C$ is the maximum over the 72 cells at least three cells from the border and responds to compact objects away from the rim. The exclusion score is $s_t = \max(G/t_G,\, C/t_C)$ with $t_G$ and $t_C$ the 99.5th percentiles of $G$ and $C$ on the calibration frames, so that $s_t > s_{\mathrm{anom}} = 1$ flags about \SI{1}{\percent} of nominal frames, and a flagged frame is withheld from the update loop by line 1 of \cref{alg:ch4-conditional-adaptation}. Nothing else happens to it: the forecast is issued as usual.

\paragraph{Synthetic faults.} No annotated sensing faults exist for any of the three sites, so detection is measured on faults rendered onto clean 2016 frames (\cref{fig:ch4-fault-mosaic}): \emph{soiling}, three to six diffuse dark patches persisting for one to three whole days; \emph{occlusion}, an opaque object of a colour and shape drawn per episode, standing in for a bird and lasting 5 to 30 minutes; \emph{banding}, horizontal gain bands and row noise redrawn every frame, modelled on an acquisition fault of the Folsom camera and lasting 1 to 5 minutes; and \emph{insect}, a body of a few pixels with legs, lasting 2 to 15 minutes. Soiling and occlusion are also rendered at the rim of the field of view, where the exclusion flag discards evidence by construction. Two further faults were used along the way and retired: sparse droplets, too small at this resolution for any detector, and a dense wet-dome rendering that was more opaque than the real frames motivating it; their rows are reported in the reproduction record and omitted here.

\paragraph{The maintenance alert.} A persistent fault calls for a different response than exclusion: someone has to clean the dome. Three alert designs were pre-registered and tested, each a causal rolling median of a frame score over the preceding 30 or 60 minutes of the same day, thresholded at its 99th percentile on 30 whole calibration days, with the criterion that every soiling episode of the corrupted stream raises an alert within two hours while alerts fire on at most \SI{1}{\percent} of fault-free frames. Their outcome is reported below in one paragraph, because all three fail, and for reasons that bear on the exclusion flag.

\subsection{Evaluation}
\label{sec:ch4-admissibility-evaluation}

\paragraph{Evaluation protocol.} The following are reported.
\begin{itemize}
  \item \textbf{Detection of synthetic faults:} for each detector, the share of faulty frames flagged at its own threshold and the AUROC against clean frames, on \num{3070} Folsom 2016 daylight frames per condition (\cref{tab:ch4-detection}); for the exclusion flag, the pre-registered criterion of at least \SI{50}{\percent} of banding and occlusion frames flagged at no more than \SI{2}{\percent} of clean frames.
  \item \textbf{Real flagged frames:} the frames the exclusion flag ranks highest on the replayed 2016 stream, one per day, inspected and grouped by the physical condition producing them (\cref{fig:ch4-flagged-frames}).
  \item \textbf{The policy under faults:} the static model, the damped policy and the damped policy with the exclusion flag active, on four streams: the clean year; a year with mixed episodes of every fault on \SI{3.9}{\percent} of daylight frames; a continuous soiling from 1 July to 15 August, then cleaned; and a permanent occluder from 1 September to the end of the year. Error on fault-free and on faulty frames, updates made on faulty frames, frames withheld, and the paired day-block bootstrap contrasts of \cref{fig:ch4-contrasts} (\cref{tab:ch4-streams}, \cref{fig:ch4-continuous}). Three further streams with mixed episodes on 1, 5 and \SI{15}{\percent} of frames were run to test whether the dose matters; it does not, and they are reduced to one sentence.
  \item \textbf{The maintenance alert:} the three designs against their pre-registered criteria.
\end{itemize}
Every arm is a single run from the checkpoint of \cref{sec:ch4-conditional-adaptation-method} on the population of \cref{tab:ch4-solar-benchmark}, and every fault stream is one draw of its schedule. The pre-registrations, the reproduction commands and the retired designs are in the repository's admissibility record.

\begin{figure}[!t]
\centering
\includegraphics[width=\textwidth]{gfx/4/ch4_admissibility_fault_mosaic.pdf}
\caption[Synthetic sensing-path faults on real frames.]{Synthetic sensing-path faults rendered onto three real Folsom 2016 frames (rows: clear, cumulus, overcast). Columns: the clean frame; soiling; occlusion, whose colour and shape are drawn per episode; the same two faults at the rim of the field of view; banding, an acquisition fault; and an insect on the dome. Every frame carries its valid irradiance label; only the image is altered.}
\label{fig:ch4-fault-mosaic}
\end{figure}

\paragraph{The embedding score does not detect sensing faults; the per-position score does.} \Cref{tab:ch4-detection} gives the result for the two detector families and the exclusion flag built from the second. At its threshold the Isolation Forest flags \SIrange{0.9}{1.7}{\percent} of frames carrying any of the six faults, indistinguishable from its \SI{1.2}{\percent} rate on clean frames, and its AUROC against clean frames is \numrange{0.54}{0.77}; the Mahalanobis and nearest-neighbour scores behave alike. The reason is the representation. The pooled embedding of an encoder trained to predict irradiance is invariant to corruptions that do not change the irradiance, which is what the invariance of its score across the columns of \cref{fig:ch4-fault-mosaic} shows. A model-free dark-pixel fraction sees opaque occlusion (\SI{22.9}{\percent}) and nothing else. The per-position score, taking the frame maximum, ranks every fault well (AUROC \numrange{0.84}{0.97}) but flags \SI{1.8}{\percent} of clean frames against the \SI{1}{\percent} it was calibrated to, because its rim cells drift between years. The exclusion flag removes the rim from the statistic and holds its calibration: it flags \SI{1.2}{\percent} of clean frames, \SI{85.4}{\percent} of banding frames and \SI{79.2}{\percent} of occlusions, passing its pre-registered criterion, together with \SI{33.0}{\percent} of insects, a few pixels at this resolution, and \SI{39.8}{\percent} of soiled frames. Its two statistics divide the work as designed: banding is caught by the global median (\SI{79.6}{\percent} alone), birds and insects by the core maximum. Faults placed at the rim are its blind spot (\SI{35.9}{\percent} of rim occlusions, \SI{8.5}{\percent} of rim soiling), the price of discarding the ring where illumination events live.

\begin{table}[!t]
\caption[Detection of synthetic sensing faults.]{Detection of synthetic sensing faults on \num{3070} Folsom 2016 daylight frames per condition. Flagged: share of frames above the detector's own threshold, calibrated on held-out 2014--2015 days to a nominal \SI{1}{\percent} false-alarm rate (\SI{0.5}{\percent} per statistic for the exclusion flag), so the clean column is the realized false-alarm rate on the evaluation year; AUROC against the clean frames. The embedding detector scores the pooled \num{2048}-dimensional embedding of the forecasting pass; the dark-pixel fraction scores raw pixels; the per-position detector scores each cell of the encoder's second stage against a per-cell memory bank, either as the frame maximum or as the exclusion flag of \cref{sec:ch4-admissibility-method}. Rim: the same fault with its centre in the outer fifth of the field-of-view radius.}
\label{tab:ch4-detection}
\centering
\small
\TableTight
\begin{tabular}{l S[table-format=2.1] S[table-format=2.1] S[table-format=2.1] S[table-format=2.1] S[table-format=2.1] S[table-format=2.1] S[table-format=2.1] S[table-format=1.2] S[table-format=1.2] S[table-format=1.2] S[table-format=1.2]}
\toprule
 & \multicolumn{7}{c}{\thead{flagged at the calibrated threshold (\si{\percent})}} & \multicolumn{4}{c}{\thead{AUROC vs. clean}}\\
\cmidrule(lr){2-8}\cmidrule(lr){9-12}
\thead{Detector} & \multicolumn{1}{c}{\thead{clean}} & \multicolumn{1}{c}{\thead{soil.}} & \multicolumn{1}{c}{\thead{soil.\\rim}} & \multicolumn{1}{c}{\thead{occl.}} & \multicolumn{1}{c}{\thead{occl.\\rim}} & \multicolumn{1}{c}{\thead{band.}} & \multicolumn{1}{c}{\thead{insect}} & \multicolumn{1}{c}{\thead{soil.}} & \multicolumn{1}{c}{\thead{occl.}} & \multicolumn{1}{c}{\thead{band.}} & \multicolumn{1}{c}{\thead{insect}}\\
\midrule
Isolation Forest (embedding) & 1.2 & 0.9 & 1.2 & 1.2 & 1.4 & 1.7 & 1.1 & 0.69 & 0.62 & 0.77 & 0.54\\
Dark-pixel fraction & 0.7 & 0.6 & 0.1 & 22.9 & 13.7 & 0.6 & 1.1 & 0.49 & 0.74 & 0.49 & 0.52\\
Per-position, frame maximum & 1.8 & 13.3 & 9.2 & 50.5 & 32.1 & 58.3 & 10.8 & 0.89 & 0.96 & 0.97 & 0.84\\
\textbf{Exclusion flag} $\max(G/t_G, C/t_C)$ & 1.2 & 39.8 & 8.5 & 79.2 & 35.9 & 85.4 & 33.0 & 0.96 & 0.99 & 0.99 & 0.94\\
\bottomrule
\end{tabular}
\end{table}

\paragraph{What the flag finds on the real stream.} \Cref{fig:ch4-flagged-frames} shows the 24 real 2016 frames the exclusion score ranks highest, one per day. The top of the ranking is sensing faults: four birds on the dome, a person standing over the camera, and three frames with the banding fault that motivated the synthetic one, together with a wet-dome frame in January and a reflection pattern inside the dome. Below them the ranking fills with rare sky, isolated cumulus, cirrus streaks, dense overcast and a red sunrise, which are under-represented in the 2014--2015 bank and are the frames behind the \SI{1.2}{\percent} clean-stream rate. This is a qualitative reading of unlabelled frames and is reported as such. It matches the synthetic result in the one respect that matters for the section: the frames the flag withholds from the policy are, at the top, unusable evidence and, further down, unusual but valid evidence.

\begin{figure}[!t]
\centering
\includegraphics[width=\textwidth]{gfx/4/ch4_admissibility_monitor_grid.pdf}
\caption[Real frames the exclusion flag ranks highest.]{The 24 real Folsom 2016 frames the exclusion score ranks highest, one per day, in rank order. Each title gives the date, the time, which statistic fired (G, the global median; C, the core maximum, whose cell is boxed) and the score as a multiple of the threshold. Birds, a person over the camera, acquisition banding, water on the dome and an internal reflection occupy the first two rows and the start of the third; unusual but valid sky fills the rest.}
\label{fig:ch4-flagged-frames}
\end{figure}

\paragraph{The policy under faults.} \Cref{tab:ch4-streams} is the substantive result of the section. It reports the static model, the damped policy and the damped policy with the exclusion flag active on the clean year and on three corrupted copies of it. Three findings follow, in order of increasing consequence.

First, learning from faulty frames does not contaminate the forecasts on clean ones. On the mixed-episode stream the policy makes \num{4423} of its updates on faulty frames, \SI{35}{\percent} more than the same frames attract when clean, because the faults raise the lagged error that triggers it; its skill on the fault-free frames is within \num{0.06} points of the clean-year run. The three dose streams say the same at 1, 5 and \SI{15}{\percent} of frames faulty (\num{-0.01}, \num{-0.05} and \num{-0.03} points), so the objection that the faults were too rare to matter does not survive. The labels of a soiled or occluded frame are valid, and an update toward a valid label from a corrupted image is a small perturbation, not a poisoning.

Second, the exclusion flag is accuracy-neutral on transient faults and cheap. On the clean year it withholds \SI{1.2}{\percent} of frames and changes skill by \num{-0.05} points ($[-0.10, -0.00]$) with MAE \SI{0.09}{\percent} lower; on the mixed-episode stream it withholds \SI{82}{\percent} of the banding frames and \SI{78}{\percent} of the occlusions in their episodes, removes a third of the updates made on faulty frames, and changes fault-free skill by the same \num{-0.05} points. It costs a third of what the embedding gate of the originating publication costs on the same runs (\num{-0.15} points, $[-0.26, -0.03]$), because what it withholds is mostly uninformative rather than mostly informative. Its justification is therefore not accuracy but budget: the updates it saves are made on frames the model cannot generalize from.

Third, and this is the finding that decides the section, withholding a persistent fault is harmful. Under continuous soiling the flag withholds \SI{38}{\percent} of the soiled frames; under the permanent occluder, \SI{87}{\percent}. The ungated policy learns through both. \Cref{fig:ch4-continuous} follows its error on the faulty frames week by week: under soiling its excess over the clean-year run on the same dates halves between the first and the last two weeks of the episode (from \num{+1.4} to \SI{+0.7}{\watt\per\square\meter}), and under the occluder its error in the last two weeks is within \SI{1}{\percent} of the clean-year run on the same dates. The gated policy, starved of the frames it would learn from, does neither: its excess under soiling grows (from \num{+2.1} to \num{+3.4}), and under the occluder it stays \SI{4}{\percent} above the clean-year run. Over the whole episode the flag costs \SI{21}{\percent} of MAE on the soiled frames ($[+11, +29]$, \num{-4.6} skill points, $[-13.4, -1.6]$) and \SI{3.2}{\percent} on the occluded ones ($[+1.9, +5.1]$, \num{-1.4} points, $[-1.9, -1.1]$). Against the static model, which cannot learn through either fault, the ungated policy is \num{13.7} skill points better on the soiled frames ($[+4.2, +43.5]$; MAE \SI{34}{\percent} lower) and \num{2.5} points better on the occluded ones ($[+1.9, +3.7]$); the gated policy keeps \num{9.1} and \num{1.1} of those points. The same holds on the mixed-dose streams, where the policy beats the static model on faulty frames by \num{4.0}, \num{14.5} and \num{7.5} points at 1, 5 and \SI{15}{\percent} of frames faulty, with the gated policy always between the two.

\begin{table}[!t]
\caption[The policy under sensing faults, with and without exclusion.]{The static model, the damped policy of \cref{eq:ch4-damped-encoder-update} and the same policy with the exclusion flag active, on the clean Folsom 2016 year and three corrupted copies of it: mixed episodes of every fault of \cref{fig:ch4-fault-mosaic} (the pre-registered stream of the exclusion flag, which also carried a dense-droplet fault later retired); one continuous soiling from 1 July to 15 August, then cleaned; and one permanent occluder from 1 September. MAE (\si{\watt\per\square\meter}) and forecast skill (\si{\percent}) on the fault-free and on the faulty frames of each stream; updates made in total and on faulty frames; share of faulty frames the flag withholds; and the paired contrast in forecast skill on the faulty frames, policy against static in the policy row and gated against ungated in the gated row, with its 95\,\% day-block bootstrap interval. Single runs.}
\label{tab:ch4-streams}
\centering
\small
\TableTight
\begin{tabular}{l S[table-format=2.2] S[table-format=2.2] c c S[table-format=6.0] c c l}
\toprule
 & \multicolumn{2}{c}{\thead{fault-free frames}} & \multicolumn{2}{c}{\thead{faulty frames}} & & & & \\
\cmidrule(lr){2-3}\cmidrule(lr){4-5}
\thead{Stream, arm} & \multicolumn{1}{c}{\thead{MAE}} & \multicolumn{1}{c}{\thead{FS}} & \thead{MAE} & \thead{FS} & \multicolumn{1}{c}{\thead{updates}} & \thead{on faulty} & \thead{withheld\\(\si{\percent})} & \thead{$\Delta$FS, faulty frames}\\
\midrule
% generated by scripts/ch4_streams_table.py from tables/ch4_admissibility_streams*.csv
\multicolumn{8}{@{}l}{\itshape clean}\\
\quad static & 31.51 & 15.68 & -- & -- & 0 & -- & -- & \\
\quad policy & 31.16 & 19.23 & -- & -- & 103435 & -- & -- & \\
\quad policy, gated by $E$ & 31.13 & 19.17 & -- & -- & 100791 & -- & 1.2 & \\
\multicolumn{8}{@{}l}{\itshape mixed episodes (3.9\,\% of frames)}\\
\quad static & 31.70 & 15.69 & 39.11 & 6.14 & 0 & -- & -- & \\
\quad policy & 31.37 & 19.31 & 29.83 & 14.56 & 105116 & 4423 & -- & +8.42 [+0.52, +20.94]\\
\quad policy, gated by $E$ & 31.32 & 19.25 & 30.11 & 14.54 & 100901 & 2938 & 30 & -0.01 [-0.29, +0.23]\\
\multicolumn{8}{@{}l}{\itshape continuous soiling (14.6\,\% of frames)}\\
\quad static & 35.63 & 15.69 & 13.10 & 7.40 & 0 & -- & -- & \\
\quad policy & 35.31 & 19.17 & 8.62 & 21.14 & 106829 & 6460 & -- & +13.74 [+4.61, +42.12]\\
\quad policy, gated by $E$ & 35.19 & 19.14 & 10.40 & 16.55 & 99985 & 2878 & 38 & -4.59 [-13.61, -1.65]\\
\multicolumn{8}{@{}l}{\itshape permanent occluder (29.7\,\% of frames)}\\
\quad static & 33.68 & 16.01 & 26.90 & 13.70 & 0 & -- & -- & \\
\quad policy & 32.97 & 19.87 & 27.05 & 16.23 & 103971 & 31143 & -- & +2.53 [+1.85, +3.68]\\
\quad policy, gated by $E$ & 32.92 & 19.82 & 27.92 & 14.80 & 74602 & 3854 & 87 & -1.43 [-1.93, -1.03]\\
\bottomrule
\end{tabular}
\end{table}

\begin{figure}[!t]
\centering
\includegraphics[width=\textwidth]{gfx/4/ch4_admissibility_continuous.pdf}
\caption[Adaptation through a persistent fault, with and without exclusion.]{Weekly MAE on the faulty frames under continuous soiling (left, 1 July to 15 August) and under the permanent occluder (right, from 1 September): the static model, the damped policy, and the same policy with the exclusion flag active; bars give the share of faulty frames the flag withholds each week. The ungated policy converges toward its clean-year error on the same dates; the gated policy, denied the frames, does not.}
\label{fig:ch4-continuous}
\end{figure}

\paragraph{The maintenance alert fails, for a reason that matters.} None of the three pre-registered alert designs meets its criteria. The frame-maximum alert raises an alert within two hours in three of four soiling episodes and fires on \SI{2.8}{\percent} of fault-free frames against the \SI{1}{\percent} bar, because a threshold calibrated on 2014--2015 flags \SI{1.5}{\percent} of 2016 frames and the rolling median amplifies clustered exceedances; restricting it to interior cells brings the false alerts to \SI{1.6}{\percent} and detects the same three episodes. The third design, a 60-minute median of the global score, separates transient from persistent events as intended, alerting during \SI{3}{\percent} of bird, banding and insect episodes, but fires on \SI{3.9}{\percent} of fault-free frames and, decisively, during only \SI{2.9}{\percent} of the frames of the continuous soiling and \SI{5.0}{\percent} of the occluded ones. Three to six soiling patches or one occluder cover a minority of the 172 cells, so the median cell score never moves; the faults the policy handles best are the ones this statistic cannot see. A per-cell persistence measure would be the next design, and is not attempted here.

\subsection{Conclusions}
\label{sec:ch4-admissibility-conclusions}
Admissibility is a property of the supervision, not of the image. The evidence of this section supports a screen that withholds transient sensing events from the update loop, reliably detected by a per-position score on an early encoder stage at a cost of one percent of clean frames and no measurable accuracy, and it argues against withholding anything that persists: a soiled or occluded frame with a valid label is evidence the policy uses, within weeks, to forecast through the fault, and a screen that removes it slows that recovery by the amount it removes. The pooled-embedding score of the originating publication detects neither class, because it measures forecasting difficulty. What remains for a screen to do is the case this section does not test: a fault in the label path, where the supervision itself is wrong and no update should be made. The design rule is stated once, in \cref{sec:ch4-summary}: withhold supervision, not images.

\paragraph{Limitations.} The evidence is limited by five constraints.
\begin{itemize}
  \item \textbf{Label-path faults are untested.} Every synthetic fault leaves the irradiance label valid. A soiled pyranometer or a clock offset corrupts the supervision, which is the case where withholding is certainly right; the section establishes only that image-path faults do not require it.
  \item \textbf{Synthetic faults only.} Detection rates and the recovery curves rest on rendered faults whose realism is not validated against annotated real ones; the real flagged frames of \cref{fig:ch4-flagged-frames} are characterized by inspection alone. The banding fault is more saturated than the real frames that motivated it, so its detection rate is optimistic.
  \item \textbf{Thresholds do not transfer across years.} Every detector was calibrated on 2014--2015 and drifts on 2016 by \numrange{0.2}{2.9} percentage points of false alarms; the exclusion flag drifts least, but none was recalibrated on the deployment year, and the same drift would be larger after relocation.
  \item \textbf{One site, one year, single runs.} All runs are Folsom 2016, one run per arm and one draw per fault schedule, so the intervals describe temporal variability only. The recovery curves of \cref{fig:ch4-continuous} are one episode each.
  \item \textbf{A guided sequence.} The three detector designs and the fault streams were pre-registered one after another, each after inspecting the previous result. Every criterion was fixed before scoring, but the sequence itself was not, and the exclusion flag is the third design tried.
\end{itemize}
The cost of the per-position detector on the target platform, the \SI{176}{\mega\byte} bank and the scoring latency, is not measured; only the second-stage activations of the forecasting pass are reused, so the added work is the nearest-neighbour search.
